
\clearpage

\section{Evaluating culture, inclusion and belonging}
\label{sec:culture}

\subsection{Provide information on how the department ensures that culture and practices support inclusion and belonging. This should include, but is not limited to, information on how the department actively considers gender equality, and EDI more broadly, in:} 
\instr{
\begin{itemize}
\item	organisation of meeting and events;
\item	images and text used in department spaces and on the department’s website; 
\item	student curricula, pedagogy, and assessment. 
\end{itemize}
}

\subsubsection{Organisation of meetings and events}\label{sec:events}
CSIT organises a variety of school meetings, committees, social events and research seminars.  
Social events are normally organised for staff at Christmas and in summer. The School follows several practices to ensure inclusion and offering opportunities for staff to be heard. 

For social events, staff are invited to register their time constraints, for example in relation to caring responsibilities. 
School meetings avoid non-core hours to ensure that everyone is able to attend.% (as discussed in \ref{sec:flex}).
School meetings are chaired by a senior member of staff other than the \ac{HOS}, who solicits agenda items from all staff a week in advance. 
During the meetings, staff are offered the opportunity to speak their minds on all school matters.
Staff meetings take place on the last Friday of the month at 15:00, and a conscious effort made to ensure that the meeting ends by 16:30. 
%The timing of the meetings is currently being discussed to ensure that the meeting ends within the core working hours (i.e. before 16:30); some topics might provoke debate among staff leading to meetings ending after 5pm. While it is important for everyone to have an opportunity to speak, the School realises that meeting exceeding core working hours are in effect negating this chance to staff with caring responsibilities.

The School organizes regular seminars to host both internal and external speakers. Recently, the School started to make specific efforts to target female speakers when issuing invitations.
All staff are invited to suggest speakers; since 2022, when suggesting a male speaker, staff are strongly encouraged to also suggest a female speaker. 
Table~\ref{tab:seminars} shows this effort has led to a better gender balance among speakers.
%the more general gender imbalance in computer science internationally impacts the female representation among seminar speakers in the School. 
%The seminar series did not take place during 2021 due to the Covid-19 pandemic and was relaunched in January 2022. 

\begin{tabbox}[9cm]{!th}

\caption{Gender distribution of speakers at CSIT seminars}
\label{tab:seminars}
\sisetup{
    group-digits=all,
    group-minimum-digits=4,
    group-separator = {,},
    table-format=4.1,
    mode=text,
    reset-text-series = false, 
    text-series-to-math = true,
    reset-text-family = false, 
    text-family-to-math = true
}
\footnotesize   
\begin{tabular}{p{3cm} 
                S[table-format=2.0]
                !{\qquad}
                S[table-format=2.0]
                !{\qquad}
                S[table-format=2.0]
                !{\qquad}
                S[table-format=2.1\%]
                }

\toprule 
Year & {Total} & {Male} & {Female} & {Female \%}\\
\midrule 
2019 & 18 & 13 & 5 & 27.8\%\\
\hdashline
2020 & 11 & 9 & 2 & 18.1\%\\
\hdashline
2021 & \multicolumn{4}{c}{No seminars due to Covid-19} \\
\hdashline
2022 & 13 & 8 & 5 & 38.4\%\\
\midrule 
Total & 42 & 30 & 12 & 28.6\%\\
\bottomrule 
\end{tabular}


\end{tabbox}

Feedback from students and the research staff focus group revealed that research staff and students lamented a lack of events and opportunities where PhD students can engage with staff and the School. 
%The research centres hosted by the School (see Section~\ref{sec:overview}) have separate communication channels to their respective members, and events are frequently not communicated to other postgraduate students. 
Most events for networking, job opportunities or professional development are organised by the various research centres hosted by the School (see Section~\ref{sec:overview}), each of which has its own dedicated communication channels such as mailing lists; while these events are open to anyone, they are not clearly communicated to those students who are not members of the research centres. 
%Students not associated to one of the centres have very limited opportunities to interact. Very few events for networking, job opportunities or professional development events are organised explicitly by the school and communicated to postgraduate students. Most are organised by the different centres with them being open to students from other centres as well, but most often this is not communicated.
This in turn weakens these researchers' sense of belonging, limiting the opportunities for a shared work culture to develop among them. %and a work culture has little opportunity to develop between all student. 
PhD students are allocated a workspace in access-controlled laboratories primarily based on the area they work in. %, rather than the centre they belong to. 
While this allows opportunities for research discussions and collaboration, it limits interaction among students working in different areas. %But since access is controlled, and generally granted only to the lab the student belongs to, collaborating and discussing with students in other fields is more difficult.

While the University offers a Research Orientation, research students also lamented a lack of orientation offered by the School.  Little information is offered to newcomers regarding the different centres, the facilities, the School, or opportunities to get involved in teaching and lab assistance.  %demonstrating and other teaching-related topics. It was felt that this does not promote an environment of inclusivity. 
%that there is little introduction given to new students by the School. 
These issues hinder a culture of inclusivity within the school, and also affect research, as it partially negates a culture of multi-disciplinary collaboration. 
Thus we propose Action~\ref{action:cult:researchinteraction} that seeks to address this.
%For this reason we propose a related action aimed at creating events and spaces where such collaborations can occur.


\begin{actionblock}{\ksspace{}Increase research interactions}{cult:researchinteraction}
\actimproveinteraction
\hfill{
    \hypersetup{hidelinks}
    \hyperlink{e2}{\gotoactionsymbol{}}
}
\end{actionblock}


%Not many policies exist for postgraduate students specifically. Most policies that concern students assume a student has responsibilities only towards courses/exams which take place in the academic year, not research. In most cases regulations/opportunities are coming from supervisor and funding agency.

\subsubsection{Images and text used in department spaces and on the department's website}
The School uses its dedicated website (\url{https://www.ucc.ie/en/compsci/}) to promote its academic programmes and research activities. A dedicated news section promotes School events, including PhD graduations and defences, undergraduate events, and outreach events, all of which include pictures of those involved.  

The School invites students from all programmes to share their thoughts on our course offerings, which are advertised on the school website. Currently, the site includes five testimonials of our undergraduate programmes, of which three of female students.

\subsubsection{Student curricula, pedagogy, and assessment}
The School's undergraduate programmes capture a diversity of interests  beyond computer science and technology: humanities, data and mathematics, psychology/human computer interaction. All programmes provide a range of pathways to cater for individual learning and interests, as shown in Table~\ref{tab:programmeoptions}.

\begin{tabbox}{!th}
\caption{Programme options}
\label{tab:programmeoptions}
\sisetup{
    group-digits=all,
    group-minimum-digits=4,
    group-separator = {,},
    table-format=4.1,
    mode=text,
    reset-text-series = false, 
    text-series-to-math = true,
    reset-text-family = false, 
    text-family-to-math = true
}
\footnotesize   
\begin{tabular}{p{5cm} 
                p{9cm}
                }

\toprule 
Programme & Options \\
\midrule 
BSc Computer Science & Computer Science, Maths and language electives\\
\hdashline
BSc Data Science and Analytics & Maths and Statistics electives\\
\hdashline
BA Psychology and Computing & Optional 9-month work placement\\
\hdashline
BA Digital Humanities and IT & Optional 9-month work placement; pursue an arts minor subject (e.g. languages, sociology, cultural studies)\\
\bottomrule 
\end{tabular}


\end{tabbox}

The School has two committees concerned with curricula, pedagogy, and assessment: the \ac{APCD} committee and the \ac{TLSE} committee.
The APCD is responsible for reviewing changes to programmes and modules delivered within the school, as well as reviewing proposals for new academic programmes. The \ac{TLSE} is focused on student experience, including drafting policies and procedures to preserve and enhance the working environment, teaching and learning experience. 

Until the start of the Athena SWAN process, the School had not actively, consciously and systematically considered how equality and diversity might enhance the curriculum, and the role these might play in pedagogy. Computer science is a broad discipline, and while certain subjects and courses are more relevant and suitable than others to discuss such issues, the Athena SWAN process has increased awareness among staff on the necessity to introduce and discuss, where relevant, gender-related elements in their teaching, for example, by using gender-neutral examples. The School plans to learn form the experience of other Computer Science schools and departments in Ireland and abroad.


\begin{actionblock}{\ksspace{}Learn from others}{cult:learnothers}
\actlearnfromothers
\hfill{
    \hypersetup{hidelinks}
    \hyperlink{e2}{\gotoactionsymbol{}}
}
\end{actionblock}


\subsection{Comment and reflect on the department’s current understanding of, and capacity to identify and address, issues and opportunities relating to equality grounds in addition to gender, as well as capacity to identify and address intersectional inequalities for staff and students.}

The School recognises intersectional inequalities for staff and students, and is particularly focused on increasing ethnic diversity, and removing barriers due to socio-economic status.
While School members come from varied ethnic backgrounds, the representation of different ethnicities varies by role, with no person in senior roles having a non-white background.
Of particular importance to the School are two socio-economic factors that, combined with gender, can significantly disadvantage staff and students. The first is income and job security inequality. Non-permanent female academic staff are over-represented (with respect to overall number in the School), as are females in below-the-bar positions (the only academic in such role is a female) (see Table~\ref{tab:acresstaff}). Insecurity of income, and lower income (at the below-the-bar grade) are socioeconomic factors that, combined with gender, produces and reinforces inequalities.
%Staff hired in the last 10 years have had significantly reduced opportunities to progress in their career due to a lack of promotion calls, and therefore to obtain higher income, which in turn reduces opportunities for securing accommodation, with new staff frequently struggling to secure adequate housing due to the current housing crisis in Ireland. 
The lack of promotion calls and opportunities for staff further ingrains this situation. %A lack of promotions means that staff are not able to secure better wages, which in turn reduces opportunities to secure accommodation.

The second factor the school is focused on is accommodation (in)security. The ongoing housing crisis in Ireland and in the Cork area affects staff and students who are renting and who may not be in a position to buy a house. Junior academic staff and research staff, especially where recently hired or on temporary contracts, frequently struggle to secure adequate accommodation. This has a disproportionate impact on female staff, as females are over-represented in all the above categories (Table~\ref{tab:acresstaff}). International staff are also both over-represented in the above categories, and disproportionately affected. Again, the intersection of gender with this socio-economic factor produces and reinforces inequalities. The School has identified the accommodation crisis as a key issue for recruitment and retention of staff at all levels, with particular concern for junior and mid-level positions. The \ac{HOS} has presented this view in higher-level university committees, up to the Academic Council, and is actively advocating for the implementation of mitigation measures aimed at reducing the impact of this crisis.
%This has disproportionately affected female members of staff, who are, for the most part, relatively recent hires. 
%The lack of job security in a number of positions, from postdoctoral researcher to Lecturer, also has socioeconomic implications that transcend and complement those of gender. Limited earning potential, combined with lack of adequate support and facilities for childcare, significantly impacts all staff hired in positions of Lecturer and above, but is specifically felt by international staff members with limited local support from extended family.

\subsection{Provide information on the department’s culture as it relates to gender equality and, where relevant, EDI more broadly, by presenting consultation findings by gender and staff category on the following areas:}
\instr{
\begin{itemize}
\item	values and traditions of the department;
\item	formal and informal structures and interactions that characterise the working and learning environment of the department, including leadership practices and behaviours;
\item	negative practices and behaviours and how these are managed by the department; 
\item	flexible working opportunities in the department;
\item	management of, and attitudes towards, family leave in the department. 
\end{itemize}
\vspace{2mm}
\noindent Where data suggests opportunity for improvement, comment and reflect. This should include reflection on any gaps between institution-level policy and practice in the department, including if the institution’s approach meets the requirements of department staff.
}

\subsubsection{Values and traditions of the department}
There is a strong collegiate culture within the School. 
All staff members are invited to social activities organised by the school. %, including monthly coffee meetings, the Christmas event, and the summer barbeque event. 
The School also facilitates various social events to get students and/or researchers together. For example, several online quiz nights were organised for postgraduate students during the pandemic. The School organises an annual reception for international students that is attended by the \ac{HOS}, Programme coordinators, and staff from the university's International Office. Furthermore, there are initiatives by individuals within the School; for example, one of the postdoctoral researchers organises weekly semi-social meeting to get researchers together, and one of the lecturers organised a well-attended PhD climbing wall event.

\subsubsection{Formal and informal structures and interactions characterising the department's working and learning environment, including leadership practices and behaviours}

All newly hired staff are assigned mentors to provide guidance on various matters in relation to their academic career and orientation within the School and University. % to ensure that they are well-integrated.

To address gender imbalance among students, the School considers numerous practices to encourage female students to join the school (see Section~\ref{sec:overview}).

Beyond gender, the \ac{HOS} recently initiated strategies to support minorities, with particular focus on ethnic minorities, present mostly in the international student cohorts. For example, the School has recently instantiated an international student mentor role to advise international students on various academic and non-academic matters. Support is offered to both international students from outside the European Union and to students from other EU countries. 
%By international students we intend here \emph{non-EU} students (as per the university definition), but such support is available to any student requesting it, with no barrier to access should a student in the \emph{EU} category request it.


The School supports female staff and students for targeted grants (see Section~\ref{sec:academic}). Similarly, the School communicates female-targeted students' scholarships (e.g., Google stipend) and encourages  students to apply for them (see Section~\ref{sec:overview}). 

%The School also encourages female students to participate or become a team member in various competitions such as the Irish programming Olympiad and the International Olympiad for Informatics.  
%The School also advertises female-targeted events (e.g., IEEE Women in Engineering) and encourage students and staff to participate in these events.  


\subsubsection{Negative practices and behaviours and their management by the department}

No specific negative behaviour by staff or students was reported in the survey and focus groups. Data accessible to the \ac{SAT} (therefore excluding any potential confidential complaint) does not point to instances of any negative behaviours. The School management will however act in accordance with University policy to address any such behaviour, should it arise in the future. 

The survey and focus groups did highlight certain negative practices in relation to family leave, inclusivity, and belonging; these are discussed in the relevant sections below. 
%For example, issues around family leave, backfill and cover, current policies requiring service length and funding for researcher discussed in Section~\ref{sec:famleave}; while issues around lack of interactions between research groups are presented in Section~\ref{sec:famleave}. 


%\subsubsection{Leadership practices and behaviours \todo{to be merged into other sections?}}


%In the following, the leadership practices discussed in one-to-one interviews to (current and former) heads and managers are presented.

%\begin{itemize}

%\item	\textbf{Staff and Student Recruitment}: 
%The school leaders recognise gender imbalance among academic staff and have taken many steps to overcome. 
%HR was consulted to ensure that job advertisements include inviting statements for females. 
%Additionally, job openings are directly communicated to qualified candidates from various institutions. 
%These actions increased the number of female applicants for recently advertised positions leading to hiring a first female lecturer in 2017, and three additional female lecturers joined the staff in 2018 and 2019. 
%It is worth noting that some positions were offered to short-listed female candidates but they declined the offer due to the high teaching load. 
% The school considers evolving opportunities to overcome the gender imbalance among academic staff. 
% For example, the school tried twice to hire female professors in HCI and AI through the \ac{SALI} scheme but that was not successful. On hiring researchers, the directors of the research centres target maintaining gender balance. A similar goal is considered when recruiting Principal and Funded Investigators for research centres. 


%\item	\textbf{Existing Staff and Students}: 

%The \ac{HOS} follows a democratic process for load assignment. 
%Staff are surveyed for their preference (teaching and service tasks), and the assignment prioritises such preference, while ensuring meeting other objectives, such as balancing gender, type -of-job, and seniority level when forming committees. The school head and administration keep close contact with year heads to ensure that student voice is heard.  



\subsubsection{Flexible Working}

The school currently supports flexible working for all staff. 
%Staff who take long-term leave (e.g., maternity or illness) are offered a reduced load upon their return. 
Most staff are allowed to work from home under the Pilot Blended Working Policy, however some staff have not been allowed due to the nature of their job.
Blended Working varies from academic to support staff due to the nature of the roles. The feedback from the PMS focus group as well as the EDI survey indicated a strong support for blended working to continue (75\% either agree or strongly agree to blended working; there is no breakdown across gender) and participants expressed a need for clarity on this.
In the PMS focus group, staff expressed near-unanimous strong support for Blended Working, which led to Action~\ref{action:cult:flexwork}. Staff acknowledged that the Blended Work policy may need tailoring to the role’s responsibilities, but that every effort should be made to seek a solution that offers flexibility to staff. Many PMS staff members highlighted the need for the University to clarify future policy. 
%Others felt strongly that flexible start/finish times should be formalised, assuming the \textit{``job can still be done.''} 


\begin{actionblock}{\ksspace{}Flexible Working}{cult:flexwork}
\actadoptflex
\hfill{
    \hypersetup{hidelinks}
    \hyperlink{e5}{\gotoactionsymbol{}}
}
\end{actionblock}
 

Working hours vary by role. Staff in focus groups said they felt comfortable discussing flexible working arrangements with their line manager. The School's reception is open daily 9:00 to 17:00 which limits some PMS staff to avail of flexible start and finish times. Most working arrangements seem to be informally arranged with line managers where applicable.  
Academics can indicate preferences for certain time slots for laboratory sessions, though  restrictions such as the availability of adequately sized rooms must also be considered. 
%For academic staff, changes to teaching schedules are facilitated where possible when asked but it should be noted that some restrictions can apply based on availability of adequately sized rooms/labs and on programmes taking certain modules. 

% Some other existing UCC policies which promote flexible working are as follows:
% \begin{itemize}[noitemsep]
% \item	Flexible Working Hours
% \item	Incentivised Shorter Working Year Scheme
% \item	Reduced Working Week
% \end{itemize}

%It is unclear from the Staff Survey or from the Focus Groups if anyone in the School is availing or has availed of the above. 

%Annual leave and when it can be taken should also be considered. 
Some restrictions to annual leave apply due to the nature of some roles; for example, academics cannot normally take annual leave during teaching periods. Similar restrictions apply to some PMS staff. 

Research students are given the option to flexible working hours, if their supervisor approves. There is also flexibility to work from home or abroad (such as the home country of international students) as long as the supervisor approves.

\subsubsection{Management and attitudes towards family Leave}%\label{sec:famleave}

Family leave policies are stipulated and enacted centrally by the university, with the School  in charge of implementing the cover requirements for staff on leave. 
During the Athena SWAN consultation process, one problematic issue was highlighted:
UCC currently maintains a requirement for staff members to have a minimum of six months of service before they are eligible for paid maternity leave. This requirement puts UCC at odds with other Irish universities, where this requirement was either never in place, or was recently removed (e.g. UCD in 2019). 
%For example, University College Dublin removed this requirement in a policy revision in February 2019. 

\begin{tabbox}[11cm]{!t}
\caption{Leave taken during 2017-2020}
\label{tab:culture:leave}
\footnotesize 

    
\begin{tabular}{p{1.4cm}
                !{\qquad}
                p{4cm}
                !{\qquad}
                p{3cm}
                R{1cm}}
    
\toprule 
Year  & Position & Type of leave & Days \\
\midrule
    
2017/18 & Research support officer & Parental & 6\\
\cdashline{2-4}
& Laboratory manager & Parental & 19 \\
\cdashline{2-4}
& Executive Assistant & Parental & 4\\
\cdashline{2-4}
& Senior research fellow & Paternity & 10 \\
\cdashline{2-4}
& Researcher & Paternity & 10\\
\hdashline
2018/19 & Research support officer & Parental & 2.5\\
\cdashline{2-4}
& Lecturer & Paid maternity & 78 \\
\hdashline
2019/20 & Lecturer & Paid maternity & 52\\ 
\cdashline{2-4}
& Lecturer & Unpaid maternity & 18\\
    
\bottomrule      
         
\end{tabular}
    
\end{tabbox}

Participants  of the academic focus group perceived a supportive environment for colleagues taking maternity and paternity leave. For example, one colleague returning from maternity leave was offered a reduction in teaching load for the first semester upon return. 

  
Participants commended both the current and previous Heads of School for their support for leave-takers. 

With respect to some forms of family-related leave, including paternity leave and parental leave, the provision of backfill is not formally required:
\begin{myquote}
    ``so it's up to the Head of School to try and find a way to cover your classes [and] reduce [your] teaching load for that semester.''
\end{myquote}
In the absence of a corresponding institutional commitment to support the arrangement of backfill/cover in these circumstances, Heads of School carry a responsibility that some focus group participants felt may be unreasonable.  It also leaves leave-takers dependent on the goodwill and understanding of their particular \ac{HOS} and puts an increased demand on the goodwill of colleagues, who might be asked to help cover work during a period of leave.  
Leave-takers can feel a sense of responsibility and guilt for burdening colleagues with increased workload; this in turn might deter people from taking statutory leave:    

% \begin{myquote}
% “I was lucky because they had the cover for me and [there was help from] other members of staff. If [other people] don't have that, who knows?” 
% \end{myquote}

\begin{myquote}
“You feel bad about [taking leave] because you're basically asking [colleagues] to [cover] [your] classes ... It's very awkward, and that has nothing to do with the willingness of the Head of School. As I said, [they’re] very supportive, but it's all managed informally. They have to try and help you without getting any [support themselves], pretty much.”  
\end{myquote}

Regarding research students, current policy states that a pregnant student should notify the supervisor and independently identify options she has given her funding. There is no specific UCC support. 
The student can take a leave of absence, but should first make sure of the consequences of this in her exact situation. There is no suitable university housing to support students with children, and most private student accommodation will not accept children.  

We address the issues described above in a single action with milestones addressing backfill, service requirements and funding for researcher respectively.


\begin{actionblock}{\ksspace{}Support Family Leave}{cult:famleave}
\actremoveimpediments
\hfill{
    \hypersetup{hidelinks}
    \hyperlink{e5}{\gotoactionsymbol{}}
}
\end{actionblock}


The cr\`{e}che services offered by the university are limited, and while both staff and students are eligible, students are given priority. Due to the limited capacity,  only children of students have been admitted for the past two years. Research students are eligible under the same rules as staff, however the cr\`{e}che is limited to the academic year, whereas research students work fulltime throughout the year.

International PhD students who move to Ireland obtain a student visa, with which bringing over their partner/family can be difficult as the visa they obtain will not allow the partner to work. This puts such families in a financially difficult position given the low level of a PhD stipend. 

